\chapter{Introducción}
\section{Motivacion}
La digitalización en las pequeñas y medianas empresas (pymes) representa un desafío crítico en la actualidad, donde la tecnología actúa como el motor fundamental para la competitividad y la eficiencia. No obstante, existe una brecha significativamente grande entre las necesidades digitales de la empresa y la realidad tecnológica de las pymes locales, las cuales arrastran una deuda técnica que limita su crecimiento y capacidad de respuesta.

El origen de este proyecto parte de la experiencia directa con una pyme cercana al autor, en la cual se han podido observar las carencias producidas por esta deuda técnica. A lo largo de un extenso periodo de observación en esta pyme y varias cercanas del sector, se ha podido comprobar cómo esta falta de digitalización deriva en un uso ineficiente de herramientas muy básicas, como hojas de cálculo, para tareas que requieren estructuras de datos más robustas y centralizadas. Esta realidad provoca la aparición de procesos manuales repetitivos, silos de información sin conectar y una dependencia excesiva de comunicaciones, como llamadas telefónicas o reuniones virtuales constantes, para cubrir la falta de un sistema de autogestión. Este problema no solo empeora la productividad de la empresa, sino que se provoca un aumento innecesario de la carga de trabajo y del estrés del personal.

En este contexto, el presente Trabajo de Fin de Grado surge de la identificación de estos patrones de ineficiencia y de la voluntad de aplicar los conocimientos en análisis de procesos y desarrollo de software adquiridos durante la carrera. El proyecto tratara de crear una solución software que busca transformar esta gestión fragmentada en un ecosistema automatizado y coherente.

En resumen, este trabajo se basa en el desarrollo de una aplicación orientada a la gestión y automatización de procesos operativos, permitiendo a las pymes optimizar su funcionamiento diario y mejorar la coordinación entre usuarios mediante herramientas de gestión visual.  Con el fin de alcanzar estos objetivos y garantizar una respuesta adaptativa a las necesidades detectadas, se utiliza una metodología ágil para la gestión de la documentación y el desarrollo de las tareas.
\section{Objetivos}

El objetivo principal de este Trabajo de Fin de Grado es el diseño e implementación de una plataforma de software centralizada que integre un motor de automatización de flujos de trabajo y herramientas de gestión visual. En particular, se pretende desarrollar una solución que permita a las pequeñas y medianas empresas (pymes) optimizar sus procesos operativos repetitivos mediante n8n, a la vez que se facilita la autogestión de tareas y recursos a través de un tablero Kanban y un sistema de planificación integrados.

El cumplimiento del objetivo principal se articula a través de los siguientes objetivos específicos:

\begin{itemize}
    \item[\textbf{OBJ-1}] \textbf{Estudio del contexto y análisis de requisitos:} Realizar un estudio sobre la deuda tecnológica en las pymes de ámbito local, identificando los procesos manuales más críticos y analizando aplicaciones de automatización e integración existentes en el mercado.
    
    \item[\textbf{OBJ-2}] \textbf{Análisis y selección de tecnologías:} Evaluar y seleccionar las tecnologías más adecuadas para la implementación de una arquitectura robusta, cconsiderando el uso de motores de automatización, bases de datos relacionales y entornos de desarrollo web modernos.
    
    \item[\textbf{OBJ-3}] \textbf{Adopción de una metodología ágil:} Aplicar un enfoque de desarrollo ágil basado en iteraciones (\textit{sprints}) que permita la gestión del cambio y la incorporación progresiva de funcionalidades.
    
    \item[\textbf{OBJ-4}] \textbf{Implementación de la solución integral:} Desarrollar el prototipo funcional aplicando prácticas sólidas de Ingeniería del Software (patrones de diseño, modularidad e interfaces de conexión) para garantizar un sistema eficiente, escalable y de alta calidad.
    
    \item[\textbf{OBJ-5}] \textbf{Evaluación y validación de la calidad:} Realizar pruebas de integración entre la aplicación centralizada y el motor de automatización para asegurar la integridad de los datos y la correcta ejecución de los flujos de trabajo diseñados.
\end{itemize}

\section{Competencias de la Intensificación}

Las competencias de la intensificación de Ingeniería del Software que se pretenden cubrir en este Trabajo de Fin de Grado son las siguientes:

\begin{itemize}
    \item[\textbf{IS-1}] \textbf{Capacidad para desarrollar, mantener y evaluar servicios y sistemas software} que satisfagan todos los requisitos del usuario y se comporten de forma fiable y eficiente, sean asequibles de desarrollar y mantener y cumplan normas de calidad, aplicando las teorías, principios, métodos y prácticas de la Ingeniería del Software.
    
    \item[\textbf{IS-3}] \textbf{Capacidad de dar solución a problemas de integración} en función de las estrategias y tecnologías disponibles.
    
    \item[\textbf{IS-4}] \textbf{Capacidad de identificar y analizar problemas y diseñar, desarrollar, implementar, verificar y documentar soluciones software} sobre la base de un conocimiento adecuado de las teorías, modelos y técnicas actuales.
    
    \item[\textbf{IS-5}] \textbf{Capacidad de identificar, evaluar y gestionar los riesgos potenciales asociados} que pudieran presentarse durante el ciclo de vida del desarrollo.
\end{itemize}
\section{Estructura del Proyecto}
En este apartado se expone la organización que seguirá el presente Trabajo de Fin de Grado, proporcionando una breve descripción del contenido de cada uno de los capítulos que componen la memoria.

\begin{itemize}
    \item \textbf{Capítulo 1. Introducción:} Se introduce el contexto del trabajo y se explica el porqué de su realización. Asimismo, se definen el objetivo principal y los objetivos secundarios, finalizando con la descripción de las competencias de la intensificación que se pretenden cubrir.
    
    \item \textbf{Capítulo 2. Antecedentes:} Se detalla el contexto del trabajo planteando un enfoque que va desde lo más general a lo particular. En este bloque se profundiza en la tecnología relevante para el sector y se realiza un estudio de aplicaciones similares existentes en el mercado.
    
    \item \textbf{Capítulo 3. Propuesta de Solución, Metodología y Herramientas:} Se plantea formalmente la idea antes de su desarrollo y se define la metodología de trabajo a seguir. Además, se enumeran las herramientas utilizadas para la gestión, documentación y desarrollo del software.
    
    \item \textbf{Capítulo 4. Descripción de la Aplicación:} Constituye el núcleo técnico del trabajo, donde se detallan los requisitos del sistema y la arquitectura diseñada. Incluye un resumen del proceso de desarrollo y la descripción funcional de la aplicación resultante.
    
    \item \textbf{Capítulo 5. Evaluación:} Se describe el proceso de validación del sistema a través de diferentes niveles, como la usabilidad y la satisfacción del usuario. Se exponen los resultados obtenidos y su correspondiente discusión.
    
    \item \textbf{Capítulo 6. Conclusiones y Trabajo Futuro:} Se resume el grado de cumplimiento de los objetivos y competencias planteados inicialmente. Para finalizar, se proponen ideas clave para la evolución futura del proyecto.
\end{itemize}
Complementariamente, la memoria incluye un apartado de \textbf{Bibliografía} con todas las referencias utilizadas y una serie de \textbf{Anexos} destinados a detallar el desarrollo del proyecto por sprints y a facilitar material adicional como el acceso a repositorios o prototipos.

